% ══════════════════════════════════════════════════════════════════════════
% NEW REMARKS — insert after \end{remark} of the Erdős–Shapiro remark
% (rem:erdos_shapiro) in chapters/03_spiral.tex
% ══════════════════════════════════════════════════════════════════════════

\begin{remark}[Visible lattice points and squarefree numbers]
\label{rem:visible_squarefree}
Two further classical interpretations of the constant $6/\pi^2$ deserve
mention, since both are geometric in nature and thus close to the spirit
of the present construction (see~\cite{Hardy2008}).

First, a lattice point $(x,y) \in \mathbb{Z}^2$ is called \emph{visible}
from the origin if no other lattice point lies on the open segment
between the origin and $(x,y)$; this holds precisely when
$\gcd(x,y)=1$. The density of visible points among all lattice points is
\begin{equation}
  \lim_{R \to \infty}
  \frac{\#\{(x,y) : \|(x,y)\| \leqslant R,\ \gcd(x,y)=1\}}
       {\#\{(x,y) : \|(x,y)\| \leqslant R\}}
  = \frac{6}{\pi^2}.
\end{equation}

Second, the density of the squarefree integers is likewise
\begin{equation}
  \lim_{x \to \infty} \frac{\#\{n \leqslant x : n \text{ squarefree}\}}{x}
  = \prod_{p} \Bigl(1 - \frac{1}{p^2}\Bigr)
  = \frac{1}{\zeta(2)} = \frac{6}{\pi^2},
\end{equation}
where the Euler product on the left is term by term the same product
that yields the spiral constant~$a$ in the present work.

The list of interpretations of the slope~$a$ thus extends to five: the
reciprocal of $\zeta(2)$, the mean coprimality density
(Remark~\ref{rem:euler_coprimality}), the density of visible lattice
points, the density of squarefree integers, and the slope of the
Archimedean spiral derived here. The visible-point interpretation is
noteworthy in the present context because it is itself a statement about
straight lines through integer points of the plane --- the same
raw material from which the spiral is built. Whether this analogy can be
made structural is an open question of this work, not an established
result.
\end{remark}

\begin{remark}[Farey fractions and the Franel--Landau criterion]
\label{rem:farey_franel}
The Farey sequence $F_N$ consists of all reduced fractions in $[0,1]$
with denominator at most~$N$. Its cardinality satisfies
\begin{equation}\label{eq:farey_count}
  |F_N| = 1 + \sum_{n \leqslant N} \varphi(n)
        \;\sim\; \frac{3}{\pi^2}\,N^2
        \;=\; \frac{a}{2}\,N^2 ,
\end{equation}
with the same coefficient $a/2$ that appears in the arc-length
asymptotics $s(\theta) \approx \tfrac{a}{2}\theta^2$ of the spiral
(Chapter~\ref{sec:bogenlaenge_spirale}). Both statements are counting
results governed by the mean of $\varphi$; the identical quadratic
form is recorded here as a formal correspondence, without a claim
that the two counts measure the same quantity.

The connection to the Riemann zeta function is provided by the
Franel--Landau criterion~\cite{Franel1924, Landau1924}: writing
$f_1 < f_2 < \cdots < f_{|F_N|}$ for the elements of $F_N$ and
$\delta_j = f_j - j/|F_N|$ for their deviations from perfect
equidistribution, the Riemann hypothesis is \emph{equivalent} to
\begin{equation}
  \sum_{j=1}^{|F_N|} |\delta_j| = O\bigl(N^{1/2+\varepsilon}\bigr)
  \qquad \text{for every } \varepsilon > 0.
\end{equation}
The quality of the equidistribution of the reduced fractions --- the
same arithmetic objects whose count carries the coefficient $a/2$ ---
therefore encodes the position of the nontrivial zeros of
$\zeta(s)$.
\end{remark}

\begin{remark}[Mertens' theorem and the Euler--Mascheroni constant]
\label{rem:mertens_gamma}
The spiral constant is the convergent second-power Euler product
$a = \prod_p (1 - p^{-2})$. Its first-power counterpart does not
converge to a positive limit; instead, Mertens'
theorem~\cite{Mertens1874} states
\begin{equation}\label{eq:mertens}
  \prod_{p \leqslant x} \Bigl(1 - \frac{1}{p}\Bigr)
  \;\sim\; \frac{e^{-\gamma}}{\ln x}
  \qquad (x \to \infty),
\end{equation}
where $\gamma = 0.5772156649\ldots$ is the Euler--Mascheroni constant.
The partial product on the left is the density of integers free of
prime factors $\leqslant x$; for $x = 3$ it equals
$(1-\tfrac12)(1-\tfrac13) = \tfrac13 = \varphi(6)/6$, the local
coprimality ratio of the congruence classes $n \equiv 1,5 \pmod 6$ on
which the sequence $k(n)$ is integer-valued.

The two products thus continue the same local datum in two different
ways: with squared factors the product converges and yields the spiral
constant~$a$; with first powers it decays logarithmically, and
$\gamma$ is the constant that calibrates this decay. The constant
$\gamma$ is tied to $\zeta$ at its pole through the Laurent expansion
$\zeta(s) = \frac{1}{s-1} + \gamma + O(s-1)$, while $a$ evaluates
$1/\zeta$ at $s = 2$; the pair $(\gamma, a)$ therefore probes the same
function at its two arithmetically distinguished points.
\end{remark}
